\hypertarget{quantum-computing-internals-with-python}{%
\chapter{Quantum Computing Internals with
Python}\label{quantum-computing-internals-with-python}}

Python is the primary language for the quantum computing revolution,
serving as the high-level interface for quantum circuit design and
hardware execution.

\hypertarget{the-quantum-stack}{%
\subsection{105.1 The Quantum Stack}\label{the-quantum-stack}}

\begin{enumerate}
\def\labelenumi{\arabic{enumi}.}
\tightlist
\item
  \textbf{High Level}: Python (Qiskit, Cirq).
\item
  \textbf{Transpiler}: Translates Python-defined circuits into
  hardware-specific gates.
\item
  \textbf{Backend}: Simulators (C++) or real QPUs (Quantum Processing
  Units).
\end{enumerate}

\hypertarget{qiskit-internals-the-quantumcircuit-object}{%
\subsection{\texorpdfstring{105.2 Qiskit Internals: The
\texttt{QuantumCircuit}
Object}{105.2 Qiskit Internals: The QuantumCircuit Object}}\label{qiskit-internals-the-quantumcircuit-object}}

A \passthrough{\lstinline!QuantumCircuit!} in Qiskit is a complex DAG
(Directed Acyclic Graph) of operations. * \textbf{Optimization}: Qiskit
uses C++ backends for circuit optimization, removing redundant gates and
mapping qubits to physical hardware topology to minimize decoherence and
gate errors.

\begin{center}\rule{0.5\linewidth}{0.5pt}\end{center}
